\section{The zero-range case of the truncated approximation}
\label{sec:hastings_zero_range}

The zero-range case of eqs.~(mid), (apb), and (facld) of
\cite[lines 285--348]{Hastings2006Solving} follows directly from commutativity.
Throughout this section the Gaussian filter is normalized as in
Definition~\ref{def:hastings_gaussian_filtered_term}, with the correction to
the printed weight in (mid) recorded in
\cite{gap:hastings05_truncated_dynamics}.

\begin{lemma}[Filtering an observable that commutes with the Hamiltonian]
    \label{thm:hastings_commuting_filter_exact}
    \lean{TNLean.Lattice.gaussianFilteredTerm_eq_of_commute,
        TNLean.Lattice.mem_groundEigenspace_apply_of_commute,
        TNLean.Lattice.starProjection_apply_eq_of_commute}
    \leanok
    \uses{def:hastings_gaussian_filtered_term}
    If $HW=WH$, then the Gaussian filter of $W$ under $H$ equals $W$ for
    every nonnegative variance. Moreover, if $H\psi=E_0\psi$, then
    $H(W\psi)=E_0W\psi$. Thus the orthogonal projection $P_0$ onto the
    $E_0$-eigenspace satisfies $P_0W\psi=W\psi$.
\end{lemma}
\begin{proof}\leanok
    Commutativity gives $e^{iHt}W=We^{iHt}$, so
    $e^{iHt}We^{-iHt}=W$. Its average against a Gaussian probability
    measure is $W$. Also $HW\psi=WH\psi=E_0W\psi$.
\end{proof}

\begin{lemma}[Commutativity of terms acting on individual sites]
    \label{thm:hastings_single_site_terms_commute}
    \lean{TNLean.Lattice.commute_of_singleton_support,
        TNLean.Lattice.commute_sum_of_singleton_support,
        TNLean.Lattice.gaussianFilteredTerm_sum_eq_of_singleton_support}
    \leanok
    Let $\Lambda$ be a finite set of sites, with a $d$-dimensional Hilbert
    space at each site. Suppose $H_i$ acts only at $i$ for every $i\in\Lambda$.
    Then $[H_i,H_j]=0$ for all $i,j$. Consequently, for every subset
    $s\subseteq\Lambda$, the sum $H_s=\sum_{j\in s}H_j$ commutes with every
    $H_i$, and the Gaussian filter of $H_i$ under $H_s$ equals $H_i$.
\end{lemma}
\begin{proof}\leanok
    \uses{thm:hastings_commuting_filter_exact}
    If $i\ne j$, the singleton supports are disjoint, so the terms commute.
    If $i=j$, the commutator also vanishes. Summing the commutation
    identities gives $[H_s,H_i]=0$, and
    Lemma~\ref{thm:hastings_commuting_filter_exact} evaluates the filter.
\end{proof}

\begin{lemma}[Exact truncation at zero interaction range]
    \label{thm:hastings_zero_range_exact}
    \lean{TNLean.Lattice.singleton_support_of_zero_range,
        TNLean.Lattice.gaussianFilteredTerm_sum_eq_of_zero_range,
        TNLean.Lattice.gaussianFilteredTerm_full_and_sum_eq_of_zero_range,
        TNLean.Lattice.starProjection_term_eq_of_zero_range,
        TNLean.Lattice.sub_starProjection_gaussianFilteredTerm_sum_eq_zero_of_zero_range}
    \leanok
    Let $\Lambda$ be a finite metric space, and suppose $H_i$ acts inside
    the closed ball of radius zero about $i$. Put $H=\sum_iH_i$ and let
    $s\subseteq\Lambda$ be arbitrary. For $q\ge0$ and $\Delta E\in\mathbb R$,
    both the full Gaussian filter $\widetilde H_i$ and the filter $M_i$
    under $H_s$ equal $H_i$. If $H\psi=E_0\psi$, then
    \begin{align*}
        \widetilde H_i & =M_i=H_i, \\
        (1-P_0)M_i\psi & =0.
    \end{align*}
    In particular, the approximation error in (apb) and the excited
    component in (facld) both vanish, without a spectral-gap assumption.
\end{lemma}
\begin{proof}\leanok
    \uses{thm:hastings_single_site_terms_commute,
        thm:hastings_commuting_filter_exact}
    A closed ball of radius zero in a metric space is the singleton of its
    centre. Apply Lemma~\ref{thm:hastings_single_site_terms_commute} to
    $s$ and $\Lambda$. The eigenspace statement in
    Lemma~\ref{thm:hastings_commuting_filter_exact} gives
    $P_0H_i\psi=H_i\psi$, hence the last identity.
\end{proof}
